% !PRINTED none
\begin{center}
\gesperrt{Siebenter Abschnitt.}

\medskip
{\Large\textbf{Theorie der Transformationsgleichungen.}}

\medskip
\rule{3cm}{0.4pt}
\end{center}

\begin{center}
\textbf{\S\,66. \gesperrt{Bildung von Transformationsgleichungen.}}
\end{center}

Nachdem nun die Teilungsgleichungen auf Transformationsgleichungen zurückgeführt sind, gehen wir an ein genaueres Studium dieser letzteren Gleichungen.

Wir nehmen $n$ ungerade an, verstehen unter $p$ die in $n$ aufgehenden Primzahlen, setzen

\weq{1}{\nu = \psi(n) = n\,\Pi\left(1+\frac{1}{p}\right),}

und bezeichnen die $\nu$ Reihen der Wurzeln der Teilungsgleichung mit $R_1, R_2, \ldots, R_\nu$.

Aus jeder dieser Reihen nehmen wir für

\weq{2}{\Omega_{\mu,\mu'} = \frac{4\mu K + 4\mu' i K'}{n}}

einen Repräsentanten $\Omega_1, \Omega_2, \ldots, \Omega_\nu$ und erhalten die $\varphi(n)$ Wurzeln einer Reihe, wenn wir in

\[
\mathrm{sn}\, m\,\Omega
\]

$m$ ein vollständiges System inkongruenter zu $n$ teilerfremder Zahlen durchlaufen lassen.

Die einfachsten Ausdrücke, die als Wurzeln von Transformationsgleichungen eingeführt werden können, sind die Produkte

% [sic] The source prints n above the product sign, although the following sentence says that h is the running index.
\weq{3}{\mathop{\Pi}\limits_{1,\,n-1}^{n} \Phi(\mathrm{sn}\,h\,\Omega),}

wenn $\Phi$ eine beliebige rationale Funktion ist, und $h$ die Reihe der Zahlen $1, 2, \ldots, n-1$ durchläuft.

Jede solche Funktion ist rational ausdrückbar durch $\mathrm{sn}\,\Omega$ und bleibt offenbar ungeändert, wenn $\Omega$ durch irgend ein $m\,\Omega$
%CONT>>
